% =============================================================================
% Chapter 1: Introduction
% =============================================================================

\chapter{Introduction}
\label{ch:intro}

Meetings are the backbone of how organizations communicate,
 yet the tools built to support them have not kept pace with how people actually speak.
In Egypt, business conversations naturally mix Egyptian Arabic and English within the same sentence.
This is not a communication problem; it is simply how professionals here talk.
The problem is that every mainstream meeting assistant on the market either struggles with Arabic entirely,
forces a fallback to Modern Standard Arabic that nobody uses in a boardroom,
or processes speech after the meeting is over. None of them can answer the question
عمر قال ايه لميار في أول الـ stand up meeting?" while the meeting is still running. This thesis builds a system that can.

\section{Problem Statement}
\label{sec:problem-statement}

We define the problem addressed in this thesis as follows. Given a live
multilingual meeting (with Egyptian Arabic--English code-switching as the
primary scenario), we must: (1)~capture audio in real time, (2)~produce a
low-latency word-by-word transcript that updates visibly to participants, (3)~attribute each spoken segment to a
distinct speaker, (4)~continuously index the emerging transcript into a
retrieval-augmented knowledge base, and (5)~answer natural-language questions
that refer to specific speakers and time windows, e.g., \textit{``هو أيمن قال ايه عن الـ budget اخر عشر دقايق؟''} The combined requirement of
low-latency streaming, code-switched accuracy, and temporal-speaker-aware
retrieval is the central technical problem.

\section{Motivation and Significance}
\label{sec:motivation}

Several real-world pain points motivate this work. Professional meeting
participants spend an estimated 30--50\% of their cognitive load on
note-taking and recall, which directly reduces their ability to engage
with the discussion. Post-meeting, the lack of a structured, speaker-
attributed transcript causes lost decisions and disputes about
responsibility. In multilingual regions such as the MENA area,
general-purpose ASR systems either fail on Egyptian Arabic phonology or
fall back to Modern Standard Arabic, neither of which matches natural
code-switched speech.

The significance of solving this problem in a single integrated system is
threefold. First, it produces measurable productivity gains for individual
users by removing the need for active note-taking. Second, it produces
measurable team-level gains by generating a structured, queryable
artifact from every meeting. Third, by targeting code-switched Egyptian
Arabic specifically, the work addresses an underserved language community
in commercial speech technology.

\section{Research Gaps, Questions, and Objectives}
\label{sec:rqs}

The literature (reviewed in \cref{ch:background}) contains strong solutions
to most individual sub-problems: large-vocabulary ASR via Whisper-family
models, streaming ASR with chunk-and-commit policies, speaker
diarization via Sortformer or similar architectures, and retrieval-
augmented generation (RAG) over vector stores. However, the integrated
problem stated in \cref{sec:problem-statement} contains a research gap:
code-switched Egyptian Arabic ASR with a low-latency stable commit policy
is not addressed by off-the-shelf systems, and temporal-speaker-aware
querying of a realtime RAG index is not natively supported by existing
meeting-assistant tools. The specific research questions addressed in this
thesis are:

\begin{enumerate}[label=\textbf{RQ\arabic*.}, itemsep=0.4em]
  \item \textbf{RQ1 (Streaming ASR for code-switched Egyptian Arabic).} To
  what extent can parameter-efficient fine-tuning (LoRA) of a Whisper-
  family model, combined with a local-agreement commit policy, produce
  realtime transcripts whose Word Error Rate (WER) on a held-out
  code-switched test set is below an operational threshold (TBD) while
  maintaining sub-second first-token latency?
  \item \textbf{RQ2 (Speaker attribution and temporal QA).} To what extent
  does combining speaker-diarization output with a temporal-aware QA agent
  allow natural-language questions that reference specific speakers and
  time windows to be answered with relevant evidence (Recall@$k$ TBD)?
  \item \textbf{RQ3 (Realtime knowledge-base freshness).} What is the end-
  to-end latency between a spoken utterance and its retrievability via the
  QA system, and what chunking / topic-windowing strategy minimizes this
  latency while preserving retrieval quality (MRR TBD)?
  \item \textbf{RQ4 (System robustness).} Does the integrated system
  remain stable (no transcript regressions, no speaker label swaps) over
  sessions of at least 60 minutes of continuous speech, and what is the
  observed failure mode distribution?
\end{enumerate}

\subsection{Objectives}

The objectives derived from the research questions are:

\begin{itemize}[itemsep=0.3em]
  \item \textbf{O1.} Curate a code-switched Egyptian Arabic--English speech
  corpus with validated transcripts and per-speaker labels.
  \item \textbf{O2.} Train and validate a LoRA-fine-tuned ASR model that
  outperforms the off-the-shelf baseline on the curated test set.
  \item \textbf{O3.} Implement a local-agreement streaming commit policy
  that produces stable word-by-word output with configurable latency.
  \item \textbf{O4.} Integrate a pretrained speaker-diarization model with
  the streaming ASR pipeline to produce speaker-attributed transcripts.
  \item \textbf{O5.} Implement a realtime RAG pipeline that ingests
  transcript chunks and supports temporal-speaker-aware retrieval.
  \item \textbf{O6.} Build and validate the full integrated system and
  report quantitative evaluation results.
\end{itemize}

\section{Proposed Solution Overview}
\label{sec:solution-overview}

The proposed solution, \meetai{}, is a realtime meeting assistant that
combines: (i)~a fine-tuned streaming ASR model for code-switched Egyptian
Arabic, (ii)~a local-agreement commit policy that finalizes words when they
appear in a configurable number of consecutive partial decodes, (iii)~a
pretrained Sortformer-based streaming speaker-diarization model, (iv)~a
realtime RAG knowledge base backed by PostgreSQL with the pgvector
extension, and (v)~a temporal-speaker-aware QA agent that parses natural-
language time and speaker references into structured retrieval queries. The
system is exposed through a thin React frontend and a polyglot
microservices backend (Node.js + Python).

\Cref{fig:high-level} shows the high-level system context. The full
component architecture is presented in \cref{ch:architecture}.

\begin{figure}[H]
\centering
\begin{tikzpicture}[
  node distance=1.2cm and 1.6cm,
  every node/.style={font=\small},
  box/.style={rectangle, draw, rounded corners, minimum width=3.2cm, minimum height=1cm, align=center, fill=blue!5},
  ext/.style={rectangle, draw, rounded corners, minimum width=3.2cm, minimum height=1cm, align=center, fill=orange!10},
  arrow/.style={-{Stealth[length=2mm]}, thick}
]
\node[ext] (user) {End user\\(browser)};
\node[box, right=of user] (fe) {React\\Frontend};
\node[box, right=of fe] (gw) {Node.js\\Gateway};
\node[box, right=of gw, yshift=0.8cm] (asr) {AI Engine\\(ASR + diarization)};
\node[box, right=of gw, yshift=-0.8cm] (qa) {QA Service\\(RAG + Gemini)};
\node[box, below=of gw] (meet) {Meetings\\Service};
\node[box, left=of meet] (auth) {Auth\\Service};
\node[box, below=2.6cm of asr] (db) {PostgreSQL\\+ pgvector};
\node[box, below=2.6cm of qa] (redis) {Redis\\(streams + tokens)};

\draw[arrow] (user) -- (fe);
\draw[arrow] (fe) -- (gw);
\draw[arrow] (gw) -- (asr);
\draw[arrow] (gw) -- (qa);
\draw[arrow] (gw) -- (meet);
\draw[arrow] (gw) -- (auth);
\draw[arrow] (asr) -- (db);
\draw[arrow] (asr) -- (redis);
\draw[arrow] (qa) -- (db);
\draw[arrow] (meet) -- (db);
\draw[arrow] (auth) -- (redis);
\end{tikzpicture}
\caption{High-level \meetai{} system context.}
\label{fig:high-level}
\end{figure}

\section{Thesis Organisation}
\label{sec:thesis-organisation}

The remainder of this thesis is organised as follows.

\begin{description}[itemsep=0.3em,leftmargin=4cm,style=nextline]
  \item[\Cref{ch:background}] reviews the relevant literature on ASR,
  speaker diarization, retrieval-augmented generation, and code-switched
  speech processing, and identifies the research gaps addressed by this
  thesis.
  \item[\Cref{ch:requirements}] presents the stakeholder analysis, the
  functional and non-functional requirements, and the data requirements
  of the proposed system.
  \item[\Cref{ch:architecture}] documents the system architecture and
  component design, including the high-level architecture, sequence
  diagrams for the key flows, the AI pipeline architecture, and the
  technology stack.
  \item[\Cref{ch:ml-methodology}] details the machine-learning
  methodology, including data preprocessing, the local-agreement
  streaming commit policy, the LoRA-based ASR fine-tuning strategy, the
  speaker-diarization integration, and the realtime RAG pipeline.
  \item[\Cref{ch:experiments}] describes the experimental setup,
  baselines, evaluation metrics, ablation studies, and statistical
  significance testing plan.
  \item[\Cref{ch:results}] presents and discusses the experimental
  results, including the model evaluation and the system evaluation, and
  discusses performance breakdown and failure modes.
  \item[\Cref{ch:conclusion}] concludes with a summary of contributions,
  critical reflection, and recommendations for future research.
  \item[\Cref{app:ethics,app:repro,app:genai}] provide the mandatory
  appendices on AI ethics, reproducibility, and generative-AI
  disclosure.
\end{description}
